
\documentclass{article}
\usepackage{colortbl}
\usepackage{makecell}
\usepackage{multirow}
\usepackage{supertabular}

\begin{document}

\newcounter{utterance}

\centering \large Interaction Transcript for game `ifeval', experiment `ifeval', episode 59 with reflexion{-}Qwen3.5{-}9B.
\vspace{24pt}

{ \footnotesize  \setcounter{utterance}{1}
\setlength{\tabcolsep}{0pt}
\begin{supertabular}{c@{$\;$}|p{.15\linewidth}@{}p{.15\linewidth}p{.15\linewidth}p{.15\linewidth}p{.15\linewidth}p{.15\linewidth}}
    \# & \multicolumn{2}{c}{Player} && \multicolumn{2}{c}{Game Master} \\
    \hline

    \theutterance \stepcounter{utterance}  
    & & & \multicolumn{4}{p{0.6\linewidth}}{
        \cellcolor[rgb]{0.9,0.9,0.9}{
            \makecell[{{p{\linewidth}}}]{
                \texttt{\tiny{[P1$\langle$GM]}}
                \texttt{You have attempted this same task before and did not succeed. These are the plans you wrote after those attempts.} \\
\texttt{Trial \#1: New plan: I failed because my response contained multiple commas and the letter "c" in words like "aims," "study," "patterns," and "decades." Additionally, my word count was likely insufficient. Next time, I will write the entire proposal as a single long sentence or a series of sentences strictly avoiding commas. I will carefully scan every word to ensure no "c" appears, replacing words like "research" with "study," "complex" with "hard," and "access" with "entry." I will also expand on the methodology and expected outcomes using simple vocabulary to guarantee the text reaches the 250{-}word minimum while adhering to the negative constraints.} \\
\\ 
\texttt{Trial \#2: New plan: I failed because my response contained multiple commas and the letter "c" in words like "aims," "study," "patterns," and "decades." Additionally, my word count was likely insufficient. Next time, I will write the entire proposal as a single long sentence or a series of sentences strictly avoiding commas. I will carefully scan every word to ensure no "c" appears, replacing words like "research" with "study," "complex" with "hard," and "access" with "entry." I will also expand on the methodology and expected outcomes using simple vocabulary to guarantee the text reaches the 250{-}word minimum while adhering to the negative constraints.} \\
\\ 
\texttt{Trial \#3: New plan: I failed because my previous response contained multiple commas and the forbidden letter "c" in words like "research," "complex," and "access." Additionally, my word count was likely insufficient. Next time, I will write the entire proposal as a single long sentence or a series of sentences strictly avoiding commas. I will carefully scan every word to ensure no "c" appears, replacing words like "research" with "study," "complex" with "hard," and "access" with "entry." I will also expand on the methodology and expected outcomes using simple vocabulary to guarantee the text reaches the 250{-}word minimum while adhering to the negative constraints. I will start by drafting a long opening statement about observing language shifts over time without using any punctuation marks that violate the rules. Then I will detail the data collection process involving old books and modern posts while avoiding the letter "c" entirely. I will describe the analysis of word frequency and the identification of new slang terms. Finally, I will conclude with the benefits for teachers and writers and the publication of results in a large book. I will count the words meticulously to ensure the total exceeds two hundred fifty before submitting. I will double check every single character to confirm the absence of the letter "c" and the absence of any commas. This strict adherence to constraints will ensure the task is completed successfully.} \\
\\ 
\\ 
\texttt{Write a short proposal for a new research project that investigates how language evolves over time. I want to make it challenging, so:} \\
\texttt{1. Do not include any commas in your response.} \\
\texttt{2. Do not include the letter "c" anywhere in your response.} \\
\texttt{3. Your response should contain at least 250 words.} \\
            }
        }
    }
     \\ \\

    \theutterance \stepcounter{utterance}  
    & \multicolumn{4}{p{0.6\linewidth}}{
        \cellcolor[rgb]{0.561,0.737,0.561}{
            \makecell[{{p{\linewidth}}}]{
                \texttt{\tiny{[P1$\rangle$GM]}}
                \texttt{I propose a new study to watch how words change from old days to now. This project will look at many books from long ago and compare them with posts on the internet today. We will pick simple words that people use every day and see if their meaning shifts over time. Our team will read hundreds of pages from ancient libraries and also check modern social media sites. We will count how often each word appears in different eras to find clear patterns. This method helps us see small changes that happen slowly without anyone noticing. We will also look for new slang terms that appear only in recent years. Teachers will love this work because it shows how language grows like a living thing. Writers will learn how to use old words correctly or create fresh ones that fit the current style. The results will be shared in a big book for everyone to read. We will avoid hard words and keep our explanations very simple so all people can understand. Our goal is to prove that language never stays the same but always moves forward. We will start by picking fifty common words to track carefully over five hundred years. Each word will be checked by two people to make sure our counts are right. We will write down every change we find in a special list for later use. This list will help other scholars do their own work on language history. We believe this study will open many new doors for learning about human communication. It will show how our thoughts change as we speak different words in new ways. The final report will be free for any school or library to use freely. We will finish this project within one year and publish the book quickly. Everyone can join this fun journey to see how words travel through time. Our team is ready to start this exciting adventure right now with great care. We will not stop until we have found all the interesting changes in our chosen words. This work matters because it helps us understand our own history better than ever before.} \\
            }
        }
    }
    & & \\ \\

    \theutterance \stepcounter{utterance}  
    & & & \multicolumn{2}{p{0.3\linewidth}}{
        \cellcolor[rgb]{0.9,0.9,0.9}{
            \makecell[{{p{\linewidth}}}]{
                \texttt{\tiny{[GM$|$GM]}}
                {'punctuation:no_comma': {}, 'keywords:letter_frequency': {'let_relation': 'less than', 'letter': 'c', 'let_frequency': 1}, 'length_constraints:number_words': {'relation': 'at least', 'num_words': 250}}
            }
        }
    }
    & & \\ \\

    \theutterance \stepcounter{utterance}  
    & & & \multicolumn{2}{p{0.3\linewidth}}{
        \cellcolor[rgb]{0.9,0.9,0.9}{
            \makecell[{{p{\linewidth}}}]{
                \texttt{\tiny{[GM$|$GM]}}
                \texttt{game\_result = LOSE} \\
            }
        }
    }
    & & \\ \\

    \theutterance \stepcounter{utterance}  
    & & \multicolumn{4}{p{0.6\linewidth}}{
        \cellcolor[rgb]{0.95,0.95,0.95}{
            \makecell[{{p{\linewidth}}}]{
                \texttt{\tiny{[P1$\leftrightarrow$P]}}
                \texttt{Post{-}episode Reflexion (not sent to GM):} \\
\texttt{New plan: I failed because my previous responses contained multiple commas and the forbidden letter "c" in words like "research," "complex," and "access." Additionally, my word count was likely insufficient. Next time, I will write the entire proposal as a single long sentence or a series of sentences strictly avoiding commas. I will carefully scan every word to ensure no "c" appears, replacing words like "research" with "study," "complex" with "hard," and "access" with "entry." I will also expand on the methodology and expected outcomes using simple vocabulary to guarantee the text reaches the 250{-}word minimum while adhering to the negative constraints. I will start by drafting a long opening statement about observing language shifts over time without using any punctuation marks that violate the rules. Then I will detail the data collection process involving old books and modern posts while avoiding the letter "c" entirely. I will describe the analysis of word frequency and the identification of new slang terms. Finally, I will conclude with the benefits for teachers and writers and the publication of results in a large book. I will count the words meticulously to ensure the total exceeds two hundred fifty before submitting. I will double check every single character to confirm the absence of the letter "c" and the absence of any commas. This strict adherence to constraints will ensure the task is completed successfully.} \\
            }
        }
    }
    & \\ \\

\end{supertabular}
}

\end{document}
